\scanpage{150}
\mspage{66}
% archive sheet number: 149

on a

\[
s\overset{(M,N)}{\boxtimes}t\simeq\varphi(s\deleted{\unsure{\times}},t),
\tag*{(4.3.1.4)}
\]

\margin{\uline{NB} On a bien sûr $(s,t)=s\overset{S^\wedge,T^\wedge}{\boxtimes}t$}

et $\varphi(S\times T)$ \deleted{est une} \inserted{$=\{s\boxtimes t\mid s\in S,t\in T\}$} est une partie strictement génératrice de $M\boxtimes N$. \deleted{Je ne \unsure{connais}}

\uline{NB} Je ne connais pas d'exemple où \deleted{\unsure{toutes|les}} les conditions envisagées ne soient satisfaites, mais doute pourtant à présent que même la plus faible le soit toujours, i.e. que $M\boxtimes N$ soit un pseudo-topos dès que $M$ et $N$ le sont.

\uline{Démonstration} \uline{de} 4.3.1. L'implication $b\Rightarrow b'$ est tautologique. L'implication $b'\Longleftrightarrow c$ résulte du

\noindent\begin{tabular}{@{}|p{.93\linewidth}@{}}
\uline{Lemme} 4.3.2. Soient $\unsure{\mathcal E|E},\deleted{\unsure{S}}X$ des \inserted{$\unsure{U|\mathcal U}$-}catégories, avec $X$ petite, $\mathcal E$ stable par petites $\underrightarrow{\lim}$,

\[
\alpha:\mathcal E\longrightarrow X^\wedge
\]

un foncteur. Pour que $\alpha$ admette un adjoint à gauche $\varphi$, il faut et il suffit que pour tout $x\in\mathrm{Ob}\,X$, le foncteur

\[
\left\{\begin{aligned}
\unsure{\zeta|\xi}&\longmapsto \mathrm{Hom}_{X^\wedge}(x,\alpha(\zeta))\deleted{\unsure{=}}=\alpha(\zeta)(x)\\
\mathcal E^{\unsure{o|\circ}}&\longrightarrow\mathrm{Ens}
\end{aligned}\right.
\]

soit représentable.
\end{tabular}

C'est évidemment nécessaire. La suffisance résulte du \deleted{\unsure{le fa}} fait que $X\subset X^\wedge$ est stricte%
